\begin{question}{27}{
    Voordat een container per vrachtwagen de Rotterdamse haven kan verlaten, wordt de container met een vaste röntgenscanner gecontroleerd op smokkelwaar.
    De vrachtwagen rijdt hierbij door een scanpoort, waarna het systeem het röntgenbeeld analyseert.
    De analysetijd $X$ (in minuten) van een willekeurige container is normaal verdeeld met verwachtingswaarde $E(X) = 6.5$ minuten en standaardafwijking $\sigma(X) = 1.23$ minuten.
}
    \subquestion{8}{
        Tijdens een nachtdienst worden $40$ containers één voor één gescand.
        De analysetijden zijn daarbij onafhankelijk van elkaar, en beginnen direct na elkaar.
        Bereken de kans dat de \textit{totale} scantijd voor deze $40$ containers minder dan $4\frac{1}{2}$ uur bedraagt.
    }
    \answer{
        Laat $X_i$ de scantijd zijn voor container $i$ (met $i = 1, 2, \dots, 40$). 
        Voor iedere individuele container geldt dat $E(X_i) = 6.5$ en $\sigma(X_i) = 1.23$ minuten. \rubric{1}
        
        We definiëren de totale scantijd $T$ voor de $40$ containers als de somvariabele
        \[
            T = X_1 + X_2 + \dots + X_{40}.\rubric{1}
        \]
        De centrale limietstelling zegt dat deze som $T$ van $n=40$ onafhankelijke normaal verdeelde variabelen ook normaal verdeeld is. \rubric{1}
        
        Voor de parameters van deze totale verwerkingstijd $T$ geldt:
        \begin{align*}
            E(T)      &= n \cdot \mu(X) = 40 \cdot 6.5 = 260 \text{ minuten},\rubric{1} \\
            \sigma(T) &= \sqrt{n} \cdot \sigma(X) = \sqrt{40} \cdot 1.23 \approx 7.7792 \text{ minuten}.\rubric{1}
        \end{align*}
        
        De gevraagde kans is $P(T < 4\frac{1}{2} \times 60) = P(T < 270)$. \rubric{1}

        Dit berekenen we als volgt:
        \begin{align*}
            P(T < 270) &= \normalcdf(\lb=-\GRinfinity, \ub=270, \mu=260, \sigma=7.7792) \\
                       &\approx 0.9007. \rubric{2}
        \end{align*}
        De kans dat de totale scantijd voor deze $40$ containers minder dan $270$ minuten bedraagt, is dus ongeveer \textbf{0.9007} (of \SI{90.07}{\percent}).
    }

    \subquestion{6}{
        Nadat een vrachtwagen door de scanpoort is gereden, rijdt de chauffeur naar de uitgang waar een geautomatiseerde slagboom staat.
        De tijd $Y$ (in minuten) die een vrachtwagenchauffeur nodig heeft om bij de slagboom aan te komen is normaal verdeeld met $\mu(Y) = 7.0$ en $\sigma(Y) = 0.91$.
        
        Als de analyse van het röntgenbeeld nog niet voltooid is op het moment dat de vrachtwagen de slagboom bereikt,
        moet de chauffeur stoppen en wachten in een speciale bufferzone tot de scan is vrijgegeven.
        
        Bereken de kans dat een willekeurige vrachtwagenchauffeur moet wachten in de bufferzone bij de slagboom.
        
        {
            \small\textbf{\underline{Hint}:} gebruik hiervoor de verschilvariabele $V = X - Y$: deze is normaal verdeeld met $\mu(V) = \mu(X) - \mu(Y)$ en $\sigma(V) = \sqrt{\sigma(X)^2 + \sigma(Y)^2}$.
        }
    }
    \answer{
        We hebben in deze situatie te maken met twee kansvariabelen $X$ en $Y$, zoals die hierboven gedefinieerd zijn.
        Een chauffeur moet wachten bij de slagboom als de analysetijd $X$ van het röntgenbeeld langer is dan de rijtijd $Y$ vanaf de scanpoort naar de slagboom: $X > Y$.\rubric{1}
        
        De kans dat de chauffeur moet wachten bij de slagboom is $P(X > Y)$, oftewel $P(V = X - Y > 0)$ berekenen.\rubric{1}

        De verschilvariabele $V = X - Y$ is normaal verdeeld met:
        \begin{align*}
            E(V)        &= \mu(X) - \mu(Y) = 6.5 - 7.0 = -0.5, \\
            \sigma(V)   &= \sqrt{\sigma(X)^2 + \sigma(Y)^2} = \sqrt{1.23^2 + 0.91^2} \approx 1.5300. \rubric{2}
        \end{align*}

        De gevraagde kans $P(V > 0)$ berekenen we dan als volgt:
        \begin{align*}
            P(V > 0)    &= \normalcdf(\lb=0, \ub=\GRinfinity, \mu=-0.5, \sigma=1.5300) \\
                        &\approx 0.3719. \rubric{2}  
        \end{align*}
        De kans dat een chauffeur moet wachten bij de slagboom bedraagt dus \textbf{0.3719} (of \SI{37.19}{\percent}).
    }

    \subquestion{6}{
        De tijd $Y$ die een willekeurige chauffeur nodig heeft om bij de slagboom aan te komen blijkt veel minder te variëren dan gedacht,
        oftewel de standaardafwijking $\sigma(Y)$ is kleiner dan gedacht.

        Neem aan dat de waarden van $\mu(X)$, $\mu(Y)$ en $\sigma(X)$ wel kloppen.
        Beredeneer, zonder berekening, of de kans dat een vrachtwagen in de bufferzone moet wachten groter dan / kleiner dan / gelijk aan de waarde berekend bij (b) is.
    }
    \answer{
        Uit de beginsituatie (subvraag b) blijkt dat het gemiddelde van de verschilvariabele $V = X - Y$ gelijk is aan $E(V) = -0.5$.\rubric{1}
        
        Omdat $E(V) < 0$, ligt het overgrote deel van de normale verdeling van $V$ links van de grens $0$. 
        De succeskans is de rechterstaart: $P(V > 0)$. \rubric{1}
        
        Wanneer de standaardafwijking $\sigma(Y)$ kleiner is, geldt dat ook voor de standaardafwijking van de verschilvariabele: $\sigma(V) = \sqrt{\sigma(X)^2 + \sigma(Y)^2}$.\rubric{1}
        
        Een kleinere standaardafwijking $\sigma(V)$ betekent dat de verdeling smaller wordt en zich compacter rondom het negatieve gemiddelde ($-0.5$) vormt.\rubric{1}
        
        Hierdoor wordt de kans $P(V > 0)$ ook steeds kleiner.\rubric{1}
        
        De kans dat een vrachtwagen bij de bufferzone moet wachten zal dus kleiner zijn dan de bij (b) berekende kans.\rubric{1}
    }

    \subquestion{7}{
        De Rotterdamse haven wil de kans verkleinen dat er filevorming plaatsvindt bij de slagboom, en wil om die reden het analyseren van de röntgenbeelden uitbesteden aan een AI-systeem.
        Dit AI-systeem moet ervoor zorgen dat de kans dat een vrachtwagen bij de slagboom moet wachten nog maar \SI{2}{\percent} is.
        Bereken hoe groot de verwachte analysetijd $\mu(X)$ maximaal mag zijn.
        Je mag aannemen dat de standaardafwijking $\sigma(X)$ gelijk blijft.

    }
    \answer{
        We hebben in deze situatie weer te maken met twee kansvariabelen $X$ en $Y$, maar nu is de verwachtingswaarde $\mu(X)$ een te bepalen onbekende.\rubric{1}
        Met andere woorden, de verschilvariabele $V = X - Y$ is nu normaal verdeeld met:
        \begin{align*}
            E(V)        &= \mu(X) - \mu(Y) = \mu(X) - 7.0 \rubric{1} \\
            \sigma(V)   &= \sqrt{\sigma(X)^2 + \sigma(Y)^2} = \sqrt{1.23^2 + 0.91^2} \approx 1.5300 
        \end{align*}

        Om de gevraagde $\mu(X)$ te berekenen, moeten we de vergelijking $P(V > 0) = 0.02$ oplossen.\rubric{1}
        Dit kunnen we doen met behulp van de \enquote{numerieke solver}:
        \begin{align*}
            E_1     &= \normalcdf(\lb=0, \ub=\GRinfinity, \mu=X-7.0, \sigma=1.5300) \\
            E_2     &= 0.02. \rubric{2}  
        \end{align*}
        Hieruit volgt een waarde van $X \approx 3.8578$.\rubric{1}
        Dit betekent dat de gemiddelde analysetijd $\mu(X)$ maximaal $3$ minuten en $52$ seconden mag zijn.\rubric{1}
    }
   
\end{question}